\documentclass[10pt,a4paper]{amsart}
\usepackage{amsmath,amssymb,amsthm,graphicx}
\usepackage{verbatim}
\newenvironment{reference}{\comment}{\endcomment}
%\newenvironment{reference}{}{}
\newenvironment{slogan}{\comment}{\endcomment}
\newenvironment{history}{\comment}{\endcomment}

% For commutative diagrams we use Xy-pic
\usepackage[all]{xy}

% We use 2cell for 2-commutative diagrams.
\xyoption{2cell}
% Theorem environments.
%
\theoremstyle{plain}
\newtheorem{theorem}[subsection]{Theorem}
\newtheorem{proposition}[subsection]{Proposition}
\newtheorem{lemma}[subsection]{Lemma}

\theoremstyle{definition}
\newtheorem{definition}[subsection]{Definition}
\newtheorem{example}[subsection]{Example}
\newtheorem{exercise}[subsection]{Exercise}
\newtheorem{situation}[subsection]{Situation}

\theoremstyle{remark}
\newtheorem{remark}[subsection]{Remark}
\newtheorem{remarks}[subsection]{Remarks}

\numberwithin{equation}{subsection}

% Macros
%
\def\lim{\mathop{\mathrm{lim}}\nolimits}
\def\colim{\mathop{\mathrm{colim}}\nolimits}
\def\Spec{\mathop{\mathrm{Spec}}}
\def\Hom{\mathop{\mathrm{Hom}}\nolimits}
\def\Ext{\mathop{\mathrm{Ext}}\nolimits}
\def\SheafHom{\mathop{\mathcal{H}\!\mathit{om}}\nolimits}
\def\SheafExt{\mathop{\mathcal{E}\!\mathit{xt}}\nolimits}
\def\Sch{\mathit{Sch}}
\def\Mor{\mathop{\mathrm{Mor}}\nolimits}
\def\Ob{\mathop{\mathrm{Ob}}\nolimits}
\def\Sh{\mathop{\mathit{Sh}}\nolimits}
\def\NL{\mathop{N\!L}\nolimits}
\def\CH{\mathop{\mathrm{CH}}\nolimits}
\def\proetale{{pro\text{-}\acute{e}tale}}
\def\etale{{\acute{e}tale}}
\def\QCoh{\mathit{QCoh}}
\def\Ker{\mathop{\mathrm{Ker}}}
\def\Im{\mathop{\mathrm{Im}}}
\def\Coker{\mathop{\mathrm{Coker}}}
\def\Coim{\mathop{\mathrm{Coim}}}

% Boxtimes
%
\DeclareMathSymbol{\boxtimes}{\mathbin}{AMSa}{"02}

%
% Macros for moduli stacks/spaces
%
\def\QCohstack{\mathcal{QC}\!\mathit{oh}}
\def\Cohstack{\mathcal{C}\!\mathit{oh}}
\def\Spacesstack{\mathcal{S}\!\mathit{paces}}
\def\Quotfunctor{\mathrm{Quot}}
\def\Hilbfunctor{\mathrm{Hilb}}
\def\Curvesstack{\mathcal{C}\!\mathit{urves}}
\def\Polarizedstack{\mathcal{P}\!\mathit{olarized}}
\def\Complexesstack{\mathcal{C}\!\mathit{omplexes}}
% \Pic is the operator that assigns to X its picard group, usage \Pic(X)
% \Picardstack_{X/B} denotes the Picard stack of X over B
% \Picardfunctor_{X/B} denotes the Picard functor of X over B
\def\Pic{\mathop{\mathrm{Pic}}\nolimits}
\def\Picardstack{\mathcal{P}\!\mathit{ic}}
\def\Picardfunctor{\mathrm{Pic}}
\def\Deformationcategory{\mathcal{D}\!\mathit{ef}}
\usepackage{cleveref}
\crefname{originalreference}{original reference}{original references}
\Crefname{originalreference}{Original reference}{Original references}
\makeatletter
\newlabel{lemma-blowup}{{0AGQ}{source}{}{sourcepointer-1}{}}
\newlabel{lemma-blowup@cref}{{[originalreference][0][]0AGQ}{[1][0][]source}{}{}{}}
\makeatother
\crefname{originallemma}{original lemma}{original lemmas}
\Crefname{originallemma}{Original lemma}{Original lemmas}
\makeatletter
\newlabel{morphisms-lemma-base-change-differentials}{{01V0}{source}{}{sourcepointer-2}{}}
\newlabel{morphisms-lemma-base-change-differentials@cref}{{[originallemma][0][]01V0}{[1][0][]source}{}{}{}}
\makeatother
\crefname{originallemma}{original lemma}{original lemmas}
\Crefname{originallemma}{Original lemma}{Original lemmas}
\makeatletter
\newlabel{morphisms-lemma-differentials-relative-immersion}{{01UZ}{source}{}{sourcepointer-3}{}}
\newlabel{morphisms-lemma-differentials-relative-immersion@cref}{{[originallemma][0][]01UZ}{[1][0][]source}{}{}{}}
\makeatother
\title{The differentials of the blowup of a regular local surface as the pushforward of the differentials of the exceptional divisor}
\author{The Stacks Project Authors}
\date{stacks-project:resolve.tex:local-snapshot-20261012; GNU FDL 1.2 or later}
\begin{document}
\maketitle
\noindent\textit{Editorial note.} All delivered mathematical text is the author's own, copied from the pinned Stacks Project chapter file resolve.tex; no statement or proof step was rewritten, reordered or removed.\par\medskip
\noindent\textit{Editorial note.} The counted proof is the author's complete argument: the conormal sequence of the closed immersion into the projective line, and the local computation with the relation defining the affine chart of the blowup.\par\medskip
\noindent\textit{Editorial note.} The label lemma-differentials-of-blowup names the counted statement, and the proof environment immediately following it in the source is delivered verbatim as its proof.\par\medskip
\noindent\textit{Editorial note.} The counted statement is carried by the author's own lemma environment in the chapter source.\par\medskip
\noindent\textit{Editorial note.} Pre-open gate: python operations/agent-authoring/tools/precard\_gate.py --source resolve.tex --label lemma-differentials-of-blowup --start 1011 --end 1042 --local-id 28109 returned RESULT=OK (L28109 free, occupancy FREE, public interval clear). Pre-assembly tag self-check: canonical tag resolved and not present in reference/public-items.jsonl.\par\medskip
\noindent\textit{Editorial note.} The statement and proof are copied from the pinned chapter file; the cited inputs are kept as external pointers with their canonical Stacks tags.\par\medskip
\noindent\textit{Allowed standard prerequisites (editorial).}\begin{itemize}
\item The blowup as a closed subscheme of the projective line with identified twist. Same chapter, the lemma embedding the blowup of the maximal ideal of a regular local surface into the projective line over the base and identifying the exceptional divisor with the special fibre. Provides the closed immersion whose conormal sequence computes the differentials of the blowup.
\item The conormal sequence for a closed immersion. Morphisms, the lemma relating the conormal sheaf of a relative closed immersion to the pullback of the differentials and the differentials of the source. Gives the exact sequence whose first arrow has to be shown to surject onto the kernel of the restriction map.
\item Differentials of the projective line restrict to the differentials of the exceptional curve. Morphisms, the lemma computing differentials under base change. Identifies the restriction of the pullback of the differentials with the differentials of the exceptional divisor over the residue field.
\end{itemize}\medskip
\ifdefined\hypertarget\hypertarget{sourcepointer-1}{}\fi
\noindent\textit{Source reference.} Original reference 0AGQ: sources/stacks-resolve/pinned-original-source.tex, source lines 666--666. Prerequisite pointer invoked at proof line(s) 1021; it names a standard fact of the same chapter that the counted argument uses, and every step of the counted proof built on it is retained verbatim in the delivered body.\par\medskip
\ifdefined\hypertarget\hypertarget{sourcepointer-2}{}\fi
\noindent\textit{Source reference.} Original lemma 01V0: sources/stacks-chow/tags.txt, source lines 2429--2429. Prerequisite pointer invoked at proof line(s) 1029; it names a standard fact of a sibling chapter that the counted argument uses, and every step of the counted proof built on it is retained verbatim in the delivered body.\par\medskip
\ifdefined\hypertarget\hypertarget{sourcepointer-3}{}\fi
\noindent\textit{Source reference.} Original lemma 01UZ: sources/stacks-chow/tags.txt, source lines 2428--2428. Prerequisite pointer invoked at proof line(s) 1027; it names a standard fact of a sibling chapter that the counted argument uses, and every step of the counted proof built on it is retained verbatim in the delivered body.\par\medskip
\begin{lemma}
\label{lemma-differentials-of-blowup}
Let $(A, \mathfrak m, \kappa)$ be a regular local ring of dimension $2$.
Let $f : X \to S = \Spec(A)$ be the blowing up of $A$ in $\mathfrak m$.
Then $\Omega_{X/S} = i_*\Omega_{E/\kappa}$, where $i : E \to X$
is the immersion of the exceptional divisor.
\end{lemma}

\begin{proof}
Writing $\mathbf{P}^1 = \mathbf{P}^1_S$, let
$r : X \to \mathbf{P}^1$ be as in Lemma \ref{lemma-blowup}.
Then we have an exact sequence
$$
\mathcal{C}_{X/\mathbf{P}^1} \to r^*\Omega_{\mathbf{P}^1/S} \to
\Omega_{X/S} \to 0
$$
see Morphisms, Lemma \ref{morphisms-lemma-differentials-relative-immersion}.
Since $\Omega_{\mathbf{P}^1/S}|_E = \Omega_{E/\kappa}$ by
Morphisms, Lemma \ref{morphisms-lemma-base-change-differentials}
it suffices to see that the first arrow defines a surjection
onto the kernel of the canonical map
$r^*\Omega_{\mathbf{P}^1/S} \to i_*\Omega_{E/\kappa}$.
This we can do locally. With notation as in the proof of
Lemma \ref{lemma-blowup} on an affine open of $X$ the morphism $f$
corresponds to the ring map
$$
A \to A[t]/(xt - y)
$$
where $x, y \in \mathfrak m$ are generators. Thus
$\text{d}(xt - y) = x\text{d}t$ and $y\text{d}t = t \cdot x \text{d}t$
which proves what we want.
\end{proof}

\end{document}
